\documentclass[11pt]{article}

\usepackage[T1]{fontenc}
\usepackage[utf8]{inputenc}
\usepackage[ngerman]{babel}
\usepackage[a4paper,margin=1in]{geometry}
\usepackage{amsmath,amssymb}
\usepackage[hidelinks]{hyperref}

\emergencystretch=4em
\sloppy

\title{Codex-Review zur arXiv-Reife des Cut-and-Project-Papers}
\author{}
\date{9. Mai 2026}

\begin{document}
\maketitle

\section*{Gesamturteil}

Nach erneutem Review des Papers und der Lean-Formalisierung ist mein Urteil:
Die Hauptformeln wirken mathematisch konsistent, der Lean-Code baut sauber,
und die empirischen Checks passen. Ich habe keine blockierende Luecke in den
Haupttheoremen gefunden.

Vor einer arXiv-Einreichung wuerde ich aber vier Punkte korrigieren bzw.
schaerfen. Zwei davon sind meines Erachtens wichtig fuer mathematische
Praezision, die anderen beiden verbessern Klarheit und Reproduzierbarkeit.

\section*{Wichtige Findings}

\begin{enumerate}
    \item \textbf{Literaturabschnitt zu Sturmian-Folgen.}
    In der Passage ueber Christoffel-Woerter und Sturmian-Sequenzen steht,
    dass die mechanische Sturmian-Folge rationaler Steigung \(p/q\) Periode
    \(p+q\) habe. In der ueblichen Codierung mit Steigung
    \(p/q \in (0,1)\) ist die Periode typischerweise \(q\), waehrend
    Christoffel-Woerter eine Blocklaenge \(p+q\) haben. Diese Stelle sollte
    unbedingt umformuliert werden, etwa indem nur die Christoffel-Blocklaenge
    \(p+q\) genannt wird oder die verwendete Steigungskonvention explizit
    gemacht wird.

    \item \textbf{Irrationaler Fall: Unbeschraenktheit genauer begruenden.}
    Die handschriftliche Argumentation nutzt die Dichte von
    \(s(\mathbb Z^2)\), um akzeptierte Punkte mit erster Koordinate
    \(m\to\pm\infty\) zu erhalten. Reine Dichte der gesamten Menge impliziert
    jedoch nicht automatisch Tail-Dichte mit vorgeschrieben grossem
    Koordinatenanteil. Der Lean-Beweis ist robuster: Er konstruiert einen
    kleinen internen Vektor mit positivem physischem Anteil und skaliert ihn.
    Fuer das Paper sollte entweder explizit Tail-Dichte der irrationalen
    Rotation zitiert werden oder die Lean-Idee kurz in den Text uebernommen
    werden.

    \item \textbf{Residuen \(c_r\): Repraesentanten fixieren.}
    \(c_r\) wird zuerst als Restklasse modulo \(D\) definiert und danach in
    Mengen der Form \(c_r+kD\) als ganzzahliger Repraesentant verwendet. Das
    ist mathematisch klar, aber fuer maximale Praezision sollte im Text
    explizit stehen: ``Waehle den Repraesentanten
    \(c_r\in\{0,\ldots,D-1\}\).''

    \item \textbf{Lean-README aktualisieren.}
    Das Paper beschreibt die neue irrationale Formalisierung korrekt, aber
    \texttt{Lean/CutAndProject/README.md} ist noch auf den frueheren Stand
    des rationalen Kerns zugeschnitten. Es nennt nur \texttt{Basic.lean},
    etwa 2.860 Zeilen, und sagt, der Entry point importiere nur
    \texttt{Basic}; tatsaechlich importiert
    \texttt{CutAndProject.lean} inzwischen auch \texttt{Irrational.lean}.
    Das sollte vor Veroeffentlichung des Companion-Repositories angepasst
    werden.
\end{enumerate}

\section*{Verifikationsstatus}

\begin{itemize}
    \item \texttt{lake build} in \texttt{Lean/CutAndProject} laeuft
    erfolgreich durch.
    \item Es wurden keine \texttt{sorry}, \texttt{admit} oder \texttt{axiom}
    Deklarationen im Lean-Code gefunden.
    \item \texttt{latexmk} meldet fuer das Paper keine Fehler; im Log stehen
    nur zwei harmlose Underfull-Box-Warnungen.
    \item \texttt{python3 src/python/verify\_paper.py} bestaetigt alle
    51.729 CSV-Testfaelle mit null Multiset-Mismatches, null
    Set-Mismatches und null negativen Sentinel-Werten.
\end{itemize}

\section*{Schlussfolgerung}

Das Paper ist in der Sache nahe an arXiv-reif. Die rationalen Multiset- und
Set-Formeln sind sauber durch die Residuenanalyse gedeckt; der geometrische
Multiplikator \(c_r\equiv-\alpha\beta^{-1}r\pmod D\) ist in Lean bis zu den
Headline-Theoremen durchgezogen. Der irrationale Fall ist inzwischen
substantiell formalisiert, inklusive lokaler Endlichkeit, bi-infiniter
Aufzaehlung und Periodenlift.

Meine Empfehlung: Vor der Einreichung die vier Findings oben einarbeiten,
insbesondere den Literaturvergleich zu Sturmian-Sequenzen und die
Unbeschraenktheitsbegruendung im irrationalen Fall. Danach sehe ich aus
diesem Review heraus keinen Grund, die arXiv-Einreichung weiter
aufzuschieben.

\end{document}
